\section{Introduction}\label{introduction}

Finite bounded physical systems resolve decisions through finitely many discrete acquisition events on bounded horizons. Under Landauer calibration \cite{landauer1961irreversibility,bennett1982thermodynamics}, exact resolution therefore carries irreducible thermodynamic cost. The formal object is a bounded decision system, represented in Lean by \texttt{Architecture}, carrying a positive degree-of-freedom count together with a canonical binary decision encoding \texttt{canonicalDP}. This counting parameter is exactly the structural rank of the encoded decision problem, it bounds the entropy of the decision quotient, and it determines the minimum per-cycle resolution cost. The quotient object is closer to zero-error and confusability-based information than to average-case coding \cite{shannon1956zero,korner1973graphs,lovasz1979shannon}. All claims are verified in Lean~4 \cite{moura2021lean4} with Mathlib support \cite{mathlib2020} and zero \texttt{sorry} placeholders.

\subsection{Central Result}

The convergence theorem proved in Section~\ref{five-way-equivalence} is:

\[
\begin{aligned}
&\underbrace{\mathrm{DOF}(A) = 1}_{\text{one-coordinate regime}}
\;\iff\;
\underbrace{\mathrm{srank}(\mathrm{canonicalDP}(A)) = 1}_{\text{structural rank}} \\
&\iff\;
\underbrace{\text{tractable sufficiency for } \mathrm{canonicalDP}(A)}_{\text{complexity}}
\;\iff\;
\underbrace{\text{minimum per-cycle thermodynamic cost}}_{\text{physics}}
\end{aligned}
\]

An imported coherence theorem gives a separate single-source reading of the same rank-$1$ point:
\[
\mathrm{SSOT}(A) \iff \mathrm{DOF}(A)=1,
\]
where $\mathrm{SSOT}(A)$ denotes the coherent single-source condition that one locus is authoritative, every remaining encoding is a derived view, and all reachable states remain coherent.

The structural-rank and thermodynamic clauses are central. The imported coherence statement is a companion interpretation of the same rank-$1$ regime, not a premise of the mathematical-physics chain.

\subsection{Contributions}

\begin{enumerate}
\item \textbf{Finite Physical Acquisition:} bounded regions admit finitely many acquisition events, those events are discrete state transitions, and exact resolution requires a sufficient coordinate set.

\item \textbf{DOF-Structural-Rank Identification:} the canonical decision problem attached to an $n$-degree-of-freedom system has structural rank $n$.

\item \textbf{Entropy-Rank Control:} the number of decision classes is at most $2^{\mathrm{DOF}(A)}$, so the decision entropy is bounded by the degree-of-freedom count.

\item \textbf{Landauer-Linear Resolution Cost:} under Landauer calibration, exact-resolution cost is bounded below by $\mathrm{DOF}(A)\,k_B T \ln 2$.

\item \textbf{Energy-Information Duality:} the same system satisfies $E \geq k_B T\,H_{\mathrm{nats}}(D)$, linking thermodynamic cost directly to the entropy of the decision quotient.

\item \textbf{Finite-Budget No-Collapse:} bounded budget, positive per-bit cost, and exponential lower-bound growth jointly preclude physical collapse of the higher-rank regime.

\item \textbf{England Replication Inequality (Theorem~\ref{thm:england}):} $\Delta S_{\min}(k) - \Delta S_{\min}(1) \geq k_B \ln k$. The proof reduces the entropy gap to counting via $k \leq 2^{k-1}$. \leanmeta{\LH{L45}}

\item \textbf{Rank-1 Convergence (Theorem~\ref{thm:five-way}):} the rank-$1$ regime is simultaneously the one-coordinate regime, the tractable sufficiency regime, and the thermodynamic ground state.
\end{enumerate}

\paragraph{Physical significance.} Once finite acquisition is fixed as the physical interface, the degree-of-freedom count becomes simultaneously an interaction dimension, an entropy bound, and a Landauer-calibrated cost coordinate. The result concerns exact resolution structure in matter.

\subsection{Scope}

The mathematical structure links structural rank, quotient entropy, complexity, and thermodynamic cost. Theorems are stated for bounded decision systems represented by the Lean object \texttt{Architecture} and their canonical decision encoding, not for arbitrary physical systems without mediation. Architectural and programming interpretations are downstream readings of the same formal core.

\subsection{Organization}

Section~\ref{foundations} defines the structural model, the finite-acquisition interface, and the canonical decision encoding. Section~\ref{probability-model} derives the structural-rank and entropy consequences of that encoding. Section~\ref{main-theorems} derives the thermodynamic consequences. Section~\ref{five-way-equivalence} proves the convergence theorem and the England inequality. Section~\ref{related-work} situates the paper relative to thermodynamics, information theory, and formalized complexity. Section~\ref{appendix-lean} describes the Lean mechanization.
